纯径向重参数化与角度相关的几何应当区分。若采用光滑可逆映射 $r=\rho(\eta)$，度量系数仍不依赖 $\theta$；在边界算子同样保持角向不变时，Fourier正交性保留隐式模态分块，只改变径向矩阵及其条件性。当前Galerkin设计每步 $O(Kd^2)$ 的独立稠密求解成本仍适用。偏心或非圆物体的映射通常在这一表示中引入角度相关系数，并使角模耦合。这是关于所选坐标与算子的结构说明，不是每一种不规则几何都必须使用某种迭代算法的定理。

若耦合隐式矩阵 $A$ 与固定预条件矩阵 $B$ 均为对称正定，且满足所假设的谱等价不等式，则精确算术PCG具有标准能量误差界\cite{r21}
\begin{equation}
 \frac{\|e_j\|_A}{\|e_0\|_A}
 \leq2\left(\frac{\sqrt{\kappa}-1}{\sqrt{\kappa}+1}\right)^j,
 \qquad \kappa=\kappa_2(B^{-1/2}AB^{-1/2})\leq\beta/\alpha.
 \label{eq:geometry-pcg-bound}
\end{equation}
对于 $\kappa>1$，这给出将相对能量误差降低到 $\varepsilon$ 所需迭代次数的充分量级 $O(\sqrt{\beta/\alpha}\log(1/\varepsilon))$，不是实测迭代数或必需下界。实际代价还包括作用 $A$、用 $B$ 求解及归约。该界不会自动适用于非对称离散、不定系统或可变/非精确预条件；仅有光滑映射也没有在本文中证明所需的 $\alpha,\beta$。伪谱度量乘积可以采用角向FFT，而已组装的稀疏耦合或短Fourier卷积可使用其他乘法。因此PCG及每次矩阵向量乘均调用FFT都不是必然要求，圆环计时无法预测这种新算子组合的成本比例。
